\documentclass[11pt,reqno]{amsart}
\usepackage[T1]{fontenc}
\usepackage{lmodern,amsmath,amssymb,amsthm,mathtools,microtype,mathrsfs}
\usepackage[margin=1.1in]{geometry}
\usepackage{longtable,booktabs,array,xr-hyper}
\usepackage[hidelinks]{hyperref}
\allowdisplaybreaks[2]
\setlength{\emergencystretch}{2em}
\numberwithin{equation}{section}
\newtheorem{theorem}{Theorem}[section]
\newtheorem{lemma}[theorem]{Lemma}
\newtheorem{proposition}[theorem]{Proposition}
\newtheorem{corollary}[theorem]{Corollary}
\theoremstyle{definition}\newtheorem{definition}[theorem]{Definition}
\theoremstyle{remark}\newtheorem{remark}[theorem]{Remark}
\newcommand{\R}{\mathbb R}\newcommand{\Z}{\mathbb Z}
\newcommand{\E}{\mathbb E}\newcommand{\dd}{\,\mathrm d}
\newcommand{\one}{\mathbf1}\newcommand{\csch}{\operatorname{csch}}
\newcommand{\arcosh}{\operatorname{arcosh}}\newcommand{\tr}{\operatorname{tr}}
\newcommand{\diag}{\operatorname{diag}}\newcommand{\supp}{\operatorname{supp}}
\newcommand{\cC}{\mathcal C}\newcommand{\cD}{\mathcal D}
\newcommand{\cP}{\mathcal P}\newcommand{\id}{\mathrm{id}}
\newcommand{\cF}{\mathcal F}\newcommand{\cK}{\mathcal K}
\newcommand{\cL}{\mathcal L}\newcommand{\dist}{\operatorname{dist}}
\raggedbottom
\hypersetup{pdftitle={Action rigidity from selective observations},pdfauthor={Qian Qi},pdfsubject={A2 revision 86: shared apparatus, singular transition and divisor-period classification}}
\title[Action rigidity from selective observations]{Action rigidity from selective observations}
\author{Qian Qi}
\date{September 18, 2026}
\subjclass[2020]{35R30, 30F30, 41A05, 53C24, 62G20}
\keywords{Action rigidity, shared nuisance, weak identification, meromorphic differentials,
finite measurements, honest confidence sets, boundary rigidity}
\begin{document}
\begin{abstract}
We study action recovery from projective matrix observations with two
unknown readout channels and selective detectors. The clock threshold
changes when the detector is shared across observation sites. With a
shared detector and constant unknown reference, two clocks determine a
nowhere locally constant action field, without supplied component labels;
one clock cannot distinguish actual nonisometric simple metrics. For
binary channels and a constant relative detector at one site, three
clocks have a sharp rank-degenerate statistical transition. We prove a
uniform inverse modulus proportional to the reciprocal product of the
two channel determinants, construct honest rank-adaptive intervals,
and give matching risk and expected-length lower bounds through an
exact marginal-preserving family. For detector primitives of meromorphic
differentials on a compact real Riemann surface, a finite divisor--period
criterion classifies regular action separation, with a sufficient clock
budget determined by the pole divisor and genus. Its genus-zero case
recovers the rational polar criterion. These results complement the
sharp pointwise polynomial inverse, finite-record geometric confidence
certificates, component coverings and exact return-lift compatibility.
Classical boundary-distance rigidity and stability supply the metric
steps, with their geometric hypotheses retained.
\end{abstract}
\maketitle
\setcounter{tocdepth}{1}
\tableofcontents
